\subsection{动态规划入门}
\ojpractice{CSP202503B}{机器人饲养指南：第37次，2025-03，第2题；定义恰好消耗v个苹果的最大收益；第三节给解答}
\subsubsection{先定义状态 再写转移}
动态规划把问题拆成可重复利用的子问题。每道题都要写出五项：状态含义、转移关系、初始值、计算顺序、答案位置。只记转移公式，容易在初值和边界处出错。

所谓无后效性，是状态已经包含决定后续所需的信息。给定状态后，不必再区分到达它的不同历史。这不表示当前选择不会影响未来，当前选择当然会改变后续状态。若过去某项信息会改变可选操作，就应把它加进状态。

例如爬楼梯，每次走1、2或3阶，令 $f[i]$ 为从0阶走到 i 阶的方案数。按最后一步分类，有 $f[i]=f[i-1]+f[i-2]+f[i-3]$。令 $f[0]=1$，负下标状态为0。这里的1表示“什么也不走”的一种空方案。按 i 从小到大计算，最终答案是 f[n]。

\subsubsection{例题 最长严格上升子序列}
\practice{B3637}{最长上升子序列}
先掌握以当前位置结尾的状态，再根据数据范围选择二次或对数优化。

输入 n 个整数，选择一些位置并保持原有顺序，使选出的值严格递增，求最大长度。$1\le n\le5000$，值的绝对值不超过 $10^9$。子序列可以不连续。第一行 n，第二行数组。
\begin{verbatim}
输入
6
3 1 2 2 5 4
输出
3
\end{verbatim}
可取 $[1,2,5]$ 或 $[1,2,4]$。设 f[i] 为\textbf{以第i个元素结尾}的最长长度，
\[f[i]=1+\max\{f[j]:j<i,\ a_j<a_i\},\]
没有合法 j 时 f[i]=1。任何以 i 结尾的较长方案，倒数第二个位置必然是一个合法 j；反过来，把 a[i] 接到该 j 的方案后面也合法。因此枚举所有前驱就覆盖全部可能。答案是所有 f[i] 的最大值，不一定是 f[n]。时间 $O(n^2)$，空间 $O(n)$。
\Code{12_lis.cpp}

\subsubsection{例题 非空最大子段和}
\practice{P1115}{最大子段和}
必须选择非空连续区间，检查全负数组。

输入 n 个整数，求一个\textbf{连续且非空}区间的最大和。$1\le n\le2\times10^5$，$|a_i|\le10^9$。
\begin{verbatim}
输入
5
-2 3 -1 4 -5
输出
6
\end{verbatim}
设 ending 为以当前位置结尾的非空区间最大和。若使用当前数 x，则最优区间要么从 x 重新开始，要么延长前一位置的最优区间：$ending=\max(x,ending+x)$。全局答案记录所有位置的最大 ending。由于要求非空，初始化为第一个元素，不能初始化为0。时间 $O(n)$，额外空间 $O(1)$。
\Code{13_maxsum.cpp}

\subsubsection{例题 数字三角形的最大路径和}
输入 n 行数字，第 i 行有 i 个数。从顶点出发，每次走到下一行的正下方或右下方，求到最后一行的最大路径和。$1\le n\le1000$，每个数的绝对值不超过 $10^6$。第一行 n，随后各行给出三角形。
\begin{verbatim}
输入
3
2
-1 3
4 -2 1
输出
6
\end{verbatim}
二维状态 f[i][j] 表示到第 i 行第 j 个位置的最优值，合法前驱只有 $(i-1,j-1)$ 与 $(i-1,j)$。初值为顶点，非法位置为负无穷。转移
\[f[i][j]=a[i][j]+\max(f[i-1][j-1],f[i-1][j]).\]
上例路径 $2,3,1$ 得6。因为每行只依赖前一行，可压缩成一维。\textbf{从右向左更新}才能确保 f[j-1] 尚未被本行覆盖。代码设置 f[1]=0 作为顶点前的虚拟起点，其他位置为负无穷。时间 $O(n^2)$，额外空间 $O(n)$。
\Code{14_triangle.cpp}

\subsubsection{二维DP与最长公共子序列}
给两个长度分别为 n、m 的序列，令 $f[i][j]$ 为两个前缀的最长公共子序列长度。若最后两个元素相等，候选为 $f[i-1][j-1]+1$；还可舍弃任意一边末元素。统一写为
\[f[i][j]=\max(f[i-1][j],f[i][j-1],\ [a_i=b_j]\text{ 时的 }f[i-1][j-1]+1).\]
空前缀长度为0，按行列递增计算。时间 $O(nm)$，滚动数组可把空间压缩到 $O(m)$，但仍不减少时间。若 n、m 都到 $10^5$，仅压缩空间仍不能通过，需要利用题目的额外性质。

\subsubsection{记忆化搜索}
记忆化按需求递归计算状态，求完后缓存。它与按顺序递推的主要区别是访问顺序。缓存键必须包含完整状态；若结果还取决于一个没有放进键里的参数，缓存会错误复用。

不要简单地以“结果为0”判断是否计算过，因为0可能是合法答案。可以单设 done 数组。一般DP要求依赖关系可以按某种顺序求值；有环图上的最短路不能靠“加个缓存”就自动解决。第11章模板展开会同时使用“正在计算”和“已完成”两种标记识别环。

\subsubsection{自测与解答}
\textbf{练习1。}数组 $[5,4,3]$ 的严格LIS长度是多少？\textbf{解答。}1。每个位置的 f 都为1；若题目要求不下降序列，应把比较条件改为不大于，而不是仍用严格小于。

\textbf{练习2。}数组 $[-5,-2,-7]$ 的非空最大子段和是多少？\textbf{解答。}-2。初始化为0会错误地选择空区间。

\textbf{练习3。}三阶楼梯、每步走1、2或3阶，共多少方案？\textbf{解答。}4种：111、12、21、3。由 $f[0]=1,f[1]=1,f[2]=2$ 得 $f[3]=4$。

\textbf{练习4。}滚动数组能否把 LCS 的 $O(nm)$ 时间降为 $O(m)$？\textbf{解答。}不能。它只减少保存的历史行数，仍需要计算 n 行各 m 个状态。
\newpage
